% FILE: 04_implementation_surface.tex
% PROJECT: ma — M&A Claim Taxonomy and Board Projection Surface

\section{Implementation Surface}
\label{sec:ma-implementation}

\subsection{Source Files}
\label{sec:ma-impl-source-files}

The \texttt{ma/} module is implemented entirely as a doctrinal and specification layer. It
contains no executable code. Its implementation surface is a corpus of 40+ Markdown
specification files organized in the following groups:

\paragraph{Master Consolidation.}
\texttt{MASTER\_m\&a-ready\_process\_intelligence\_framework.md} (598 lines, 31,250 bytes)
consolidates all nine claim taxonomies with full mathematical formulas, the Four-Pillar
Admissibility Protocol, the Five-Step Auditor Evidence Path, the Cryptographic Receipt Schema,
the Synergy Proof Models, and the Operational Debt Taxonomy. It is the single authoritative
reference for the complete framework.

\paragraph{Claim Taxonomy Definitions.}
Nine \texttt{define\_*.md} files, one per claim taxonomy:
\texttt{define\_board\_claim\_taxonomy.md}, \texttt{define\_diligence\_claim\_taxonomy.md},
\texttt{define\_synergy\_claim\_taxonomy.md}, \texttt{define\_control\_claim\_taxonomy.md},
\texttt{define\_process\_asset\_claim\_taxonomy.md},
\texttt{define\_process\_liability\_claim\_taxonomy.md},
\texttt{define\_integration\_risk\_taxonomy.md},
\texttt{define\_operational\_debt\_taxonomy.md},
\texttt{define\_scalability\_claim\_taxonomy.md}.

\paragraph{Validation Framework Definitions.}
Four files defining the stakeholder evidence requirements:
\texttt{define\_buyer\_reliance\_requirements.md} (replication bounds, VDR completeness rules),
\texttt{define\_seller\_defensibility\_requirements.md} (behavioral profiles, process drift
audits), \texttt{define\_board-admissible\_claim\_requirements.md} (Four-Pillar Protocol),
\texttt{define\_auditor\_evidence\_path.md} (Five-Step Protocol).

\paragraph{Slide-to-Evidence Mapping Files.}
Four mapping specifications:
\texttt{define\_slide-to-receipt\_map.md} (cryptographic JSON receipt schema),
\texttt{define\_slide-to-replay\_map.md} (token game replay, optimal alignments),
\texttt{define\_slide-to-residual\_map.md} (residual weight and entropy),
\texttt{define\_slide-to-public-standard\_map.md} (XES, OCEL 2.0, BPMN, POWL conformance).

\paragraph{Strategic Frameworks.}
\texttt{REVERSE\_PORTER\_FIVE.md} (process evidence as structural barrier to entry),
\texttt{hostile-takeover-receipts.md} (v30.1.1 HTR taxonomy),
\texttt{synergy-proof-claims-v30.md} (Process Equivalence Model, Execution Path Intersection,
Capability Subsumption Theorem),
\texttt{ACQUISITION\_READY\_PROCESS\_INTELLIGENCE.md},
\texttt{EXECUTIVE\_BRIEF\_\_acquisition-ready-process-intelligence.md}.

\paragraph{Renderer Specification.}
\texttt{M\&A\_BOARD\_PROJECTION\_RENDERER\_SPECIFICATION.md} (25,258 bytes) defines the
seven-step pipeline from claims JSON to sealed PowerPoint and Excel artifacts.

\subsection{Commands}
\label{sec:ma-impl-commands}

The \texttt{ma/} module defines no CLI commands directly. It specifies the contract for
commands that the Board Projection Renderer must expose. Per the Renderer Specification, the
invocation sequence is:

\begin{enumerate}
  \item Populate claims JSON array from transaction data room.
  \item Execute \texttt{wasm4pm} queries to produce receipt JSON files.
  \item Invoke \texttt{Tera2} template engine against claims JSON and receipt files.
  \item Invoke \texttt{pptx-rs} to assemble the PowerPoint deck.
  \item Invoke \texttt{calamine} to assemble the Excel workbook.
  \item Execute validation protocol: verify all receipts and witness markers.
  \item Seal receipt ledger with cryptographic signatures.
\end{enumerate}

The Readiness Report (\texttt{ma-deck-rendering-readiness.md}, status: AUTHORITY STACK
COMPLETE, 2026-06-01) confirms that steps 3--7 are fully specified and await only the
claims data (step 1) and receipt generation (step 2) from a live transaction data room.

\subsection{Data and Ontology Surfaces}
\label{sec:ma-impl-data}

\paragraph{Event Log Standards.}
The module mandates XES (IEEE 1849-2016) or OCEL 2.0 as the only admissible event log
formats. XES provides case/activity/timestamp structure for single-object processes; OCEL 2.0
provides object-type/event/relationship structure for multi-object processes (Order-to-Cash,
Procure-to-Pay). All logs must include W3C PROV-O provenance metadata mapping events to
source system WAL sequence numbers.

\paragraph{Process Model Formats.}
Petri nets (WF-nets), process trees, POWL structures, and BPMN with formal semantics are all
admissible. All submitted models must carry a soundness certificate. For financial processes,
$k$-boundedness $k \le 5$ is required.

\paragraph{Receipt Schema.}
The JSON receipt schema follows RFC 8785 JCS canonical form. Required fields: \texttt{slide\_id}
(UUID), \texttt{slide\_title}, \texttt{assertion\_text}, \texttt{target\_log\_hash} (BLAKE3),
\texttt{process\_model\_hash} (BLAKE3), \texttt{query\_definition} (engine, query URI,
parameters), \texttt{verification\_results} (fitness, precision, throughput, financial
impact), \texttt{validator\_signature} (Ed25519).

\paragraph{VDR Structure.}
The Virtual Data Room is organized as six subdirectories:
\texttt{/event-logs/} (XES/OCEL 2.0 files),
\texttt{/models/} (Petri nets, process trees, POWL),
\texttt{/receipts/} (JSON receipt files),
\texttt{/conformance/} (alignment reports),
\texttt{/provenance/} (PROV-O lineage files),
\texttt{/synergy/} (behavioral similarity reports, bisimulation proofs).

\subsection{Templates and Rendered Artifacts}
\label{sec:ma-impl-templates}

\paragraph{Board Presentation Template.}
The Renderer Specification defines a \texttt{Tera2}-templated PowerPoint deck structure. Each
slide receives: a claim identifier, the assertion text, the supporting formula, the numeric
result (fitness, precision, financial impact), and a receipt UUID watermark. The watermark
enables auditors to locate the corresponding VDR receipt file by UUID.

\paragraph{Excel Workbook.}
The \texttt{calamine}-assembled workbook provides tabular representations of all claim
formulas, input parameters, computed outputs, and receipt links. It enables buyers to inspect
the arithmetic behind each board slide independently of the presentation layer.

\paragraph{Receipt Ledger.}
The sealed receipt ledger is a JSON Lines file containing all receipt records, ordered by
slide sequence, with a final ledger root hash (BLAKE3 over all receipt hashes in order) and
an Ed25519 signature over the root hash. The ledger root hash is embedded in the PowerPoint
file metadata, binding the presentation to its evidentiary substrate.
